\chapter*{Введение}\label{ux432ux432ux435ux434ux435ux43dux438ux435}
\addcontentsline{toc}{chapter}{Введение}

В индустрии видеоигр и сельскохозяйственных симуляторов постоянно возникает потребность в большом количестве названий: сортов растений, пород животных, продуктов и географических объектов. Ручное создание такого объёма наименований требует времени и нередко приводит к повторам. Автоматическая генерация позволяет получать новые варианты с заданными свойствами и единым стилем.

Актуальность темы связана с распространением генеративных моделей в обработке естественного языка и потребностью в инструментах нейминга для аграрной и игровой сфер. Для применения LSTM и трансформеров к названиям растений и продуктов необходимо адаптировать архитектуру, подготовить предметный датасет и определить критерии оценки. Это требует понимания принципов работы моделей и их ограничений.

Целью курсового проекта является разработка и исследование методов машинного обучения для генерации названий растений и аграрных объектов виртуальных миров, а также создание программного приложения, реализующего генеративный пайплайн и обеспечивающего удобный пользовательский интерфейс.

Для достижения поставленной цели необходимо решить следующие задачи:

Провести теоретический анализ генеративных моделей машинного обучения, применяемых для генерации текстовых последовательностей, включая рекуррентные нейронные сети (LSTM) и трансформеры.

Изучить методы оценки качества сгенерированных названий: перплексию, стилевое соответствие, новизну и разнообразие.

Обосновать выбор инструментов и разработать архитектуру программного приложения для генерации названий растений и продуктов.

Сформировать и подготовить датасет названий растений и продуктов, включающий разметку по категориям и длине названия.

Реализовать на языке Python генеративную LSTM-модель, классификатор стилей и модули скоринга, а также создать пользовательский интерфейс на базе Gradio.

Провести экспериментальное исследование качества обученных моделей, сравнить результаты с базовыми методами (марковские цепи, n-граммы) и проанализировать полученные данные.

Сформулировать выводы о целесообразности применения выбранных методов для задачи генерации названий в аграрной предметной области.

Объект исследования -- методы машинного обучения для генерации текстовых последовательностей; предмет -- архитектура, обучение и оценка генеративных моделей. Методы: анализ литературы, сравнение архитектур, проектирование и экспериментальное тестирование.
